\clearpage
\VerdierPhysicalPageAnchor{0202}%
\begingroup\small\setlength{\parskip}{0pt}\setlength{\emergencystretch}{1.5em}
\setlength{\abovedisplayskip}{4pt plus 2pt minus 2pt}%
\setlength{\belowdisplayskip}{4pt plus 2pt minus 2pt}%

\VerdierSemanticAnchor{ch3d-2-3-3-proof-q-r-step}
\[
\alpha\colon L^{-r+1}\longrightarrow Z^{-r+1},
\qquad\text{其中 }L^{-r+1}\text{ 是 }M\text{ 的对象}.
\]
令 $X_{q,r}=X_{q,r-1}(\alpha,-r+1)$ 且
$v_{q,r}=p(\alpha,-r+1)$（2.1.1）。显然 $X_{q,r}$ 与 $v_{q,r}$ 具有性质
$I')_r$、$III')_r$、$IV')_r$。此外，由复形 $X_{q,r}$ 的定义
（2.1.1），有正合列
\[
0\longrightarrow (X_{q,r})^{-r}\longrightarrow
L^{-r}\oplus L^{-r+1}\longrightarrow Z^{-r+1}\longrightarrow0.
\]
由性质 $(E_4)$，$(X_{q,r})^{-r}$ 属于 $M'$；由（2.2.6），
\[
M\text{-}\dim((X_{q,r})^{-r})
=M\text{-}\dim(Z^{-r+1})-1
=M\text{-}\dim(Y^q)-r-q,
\]
这正是性质 $II')_r$。

\VerdierSemanticAnchor{ch3d-2-3-3-proof-q-completion}
因此可以继续构造。令 $X_{q,-q+M\text{-}\dim(Y^q)}=X_q$，并把复合
$v_{q,l}$ 所得的态射记为 $u_q\colon X_q\longrightarrow X_{q-1}$。
显然 $X_q$ 与 $u_q$ 具有性质 $1)_q$、$2)_q$、$3)_q$、$4)_q$、
$5)_q$。命题的第一项断言得证。为证明第二项断言，只需注意在附加假设下，
对适当的 $n_0$，可在上述构造中取 $X_{n_0}=Y$。

\VerdierSemanticAnchor{ch3d-2-3-4}
\medskip\noindent\textbf{注 2.3.4.}\quad
在复形 $X_{n_0}$ 的构造中（参见命题（2.3.3）的证明），我们实际上证明了
如下结果。设 $M\subset M'$ 是 $\mathcal A$ 的两个对象集合，并且
$0\in M$，集合 $M'$ 具有性质 $(E_1)$ 与 $(E_4)$，而 $M'$ 的每个对象
都是 $M$ 中某个对象的商对象。设
$Y\in\operatorname{Ob}\operatorname{comp}^{-}(\mathcal A)$ 是各分量均属于
$M'$ 的复形。则存在各分量均属于 $M$ 的复形
$X\in\operatorname{Ob}\operatorname{comp}^{-}(\mathcal A)$，以及拟同构
$u\colon X\longrightarrow Y$。

\VerdierSemanticAnchor{ch3d-2-3-5}
\medskip\noindent\textbf{2.3.5.}\quad
把本节结果应用于 $\mathcal A$ 的反范畴所得到的结论，留给读者陈述。若
$\mathcal A$ 的对象集合 $M$ 在被视为 $\mathcal A^o$ 的对象集合时，集合
$M$ 具有（2.2.1）中的性质 $(E_i)$，则称它具有性质 $(E_i)^o$，其中
$1\leq i\leq6$。

\VerdierSemanticAnchor{ch3d-2-4}
\subsection*{2.4. 应用。}

\VerdierSemanticAnchor{ch3d-2-4-1}
\noindent\textbf{命题 2.4.1.}\quad
\textit{a) 设 $\mathcal A$ 是阿贝尔范畴，$M$ 是 $\mathcal A$ 的对象集合，
并具有性质 $(E_1)$（2.2.1）。范畴 $K^{-}(M)$（1.1.4）关于
$\operatorname{Ac}(\mathcal A)$ 是左局部化的（第二章，2.3.5），当且仅当集合
$M$ 具有性质 $(E_2)$ 与 $(E_3)$。}

\noindent\textit{b) 设 $M$ 是具有性质 $(E_1)$、$(E_2)$、$(E_3)$ 的对象集合，
并设 $M^{\wedge}$ 是（2.2.1）中定义的与 $M$ 关联的集合。自然函子}
\[
K^{-}(M)/\operatorname{Ac}^{-}(M)
\longrightarrow
K^{-}(M^{\wedge})/\operatorname{Ac}^{-}(M^{\wedge})
\]
\endgroup
